\section{Introduction}
Cine magnetic resonance imaging (MRI) is the reference for measuring the volumes, function and mass of the heart \citep{schulzmenger2013}. Once the left ventricle (LV), the right ventricle (RV) and the myocardium are segmented at end-diastole (ED) and end-systole (ES), the diagnosis of the main cardiomyopathies largely follows from a few quantities: ejection fractions, indexed volumes, myocardial mass and wall thickness \citep{pinto2016,elliott2014,marcus2010}. The ACDC challenge \citep{bernard2018} made this task a benchmark: from segmentations, classify patients into five groups of equal size. The best entries reached 92 to 96\% accuracy on its test set with a few dozen handcrafted features and standard classifiers \citep{isensee2017,khened2017,wolterink2017,cetin2017}.

We work on a course version of this challenge: 100 labeled training patients, 50 test patients whose LV cavity is not given, and test labels hidden behind a leaderboard. With 100 patients, one error is one point of accuracy and five points on a fold. Accuracies can then only be trusted if the model selection never sees the patients used for scoring \citep{varma2006,cawley2010}, and differences of one or two points between models are within the noise. We therefore built the whole pipeline as scripts, with every accuracy estimated by nested cross-validation repeated over several splits.

\paragraph{Contributions}
(1) A study of how much single-split model selection inflates accuracies on this data, compared with nested cross-validation (Section~\ref{sec:protocol-study}).
(2) An evaluation, on the training set, of the reconstruction of the missing LV cavity, with a fallback for slices where the myocardium is not closed (Section~\ref{sec:lv}).
(3) A check of the consistency between training features computed with the true LV and test features computed with a reconstructed LV (Section~\ref{sec:protocol}).
(4) Myocardial wall-thickness features and a benchmark of 10 classifiers on 3 feature sets under the same protocol (Section~\ref{sec:bench}).
(5) An analysis of the final model, its errors and its test predictions (Section~\ref{sec:final}).
The appendix gives the full tables and starts with its own table of contents.
